% ============================================================
% Appendix
% CT-MUSIQ Undergraduate Thesis
% ============================================================

\chapter{Supplementary Material}
\label{app:supplementary}

This appendix provides supplementary technical material supporting the main thesis,
including the full training log for the primary CT-MUSIQ run, detailed pseudocode for
the training loop, an extended discussion of the hash-based positional encoding, and
implementation notes for reproducing all experiments.

\section{Full Training Log Summary}
\label{app:training_log}

Table~\ref{tab:full_training_log} presents epoch-by-epoch training log data for the
primary CT-MUSIQ 100-epoch run. The table shows the learning rate, training loss,
validation PLCC, SROCC, KROCC, and Aggregate at each selected epoch. The full log is
available in \texttt{results/ct\_musiq/ct\_musiq\_training\_log.csv}.

\begin{table}[h]
    \centering
    \caption{Selected epochs from the CT-MUSIQ primary training log.
             Stage 1 (epochs 1--5): frozen encoder. Stage 2 (epochs 6--100): full fine-tuning.
             Best val.\ aggregate checkpoint is highlighted.}
    \label{tab:full_training_log}
    \small
    \begin{tabular}{rrrrrrr}
        \toprule
        \textbf{Epoch} & \textbf{Stage} & \textbf{LR} & \textbf{Train Loss}
        & \textbf{PLCC} & \textbf{SROCC} & \textbf{Aggregate} \\
        \midrule
        1   & 1 & $3.0 \times 10^{-4}$ & 0.877 & 0.877 & 0.901 & 2.513 \\
        2   & 1 & $3.0 \times 10^{-4}$ & 0.432 & 0.881 & 0.917 & 2.559 \\
        3   & 1 & $3.0 \times 10^{-4}$ & 0.360 & 0.884 & 0.906 & 2.534 \\
        4   & 1 & $3.0 \times 10^{-4}$ & 0.315 & 0.951 & 0.953 & 2.733 \\
        5   & 1 & $3.0 \times 10^{-4}$ & 0.343 & 0.886 & 0.926 & 2.591 \\
        \midrule
        6   & 2 & $5.0 \times 10^{-5}$ & 0.346 & 0.947 & 0.949 & 2.714 \\
        7   & 2 & $5.0 \times 10^{-5}$ & 0.269 & 0.939 & 0.947 & 2.695 \\
        8   & 2 & $5.0 \times 10^{-5}$ & 0.281 & 0.940 & 0.942 & 2.688 \\
        9   & 2 & $\sim 4.8 \times 10^{-5}$ & 0.272 & 0.957 & 0.962 & 2.765 \\
        10  & 2 & $\sim 4.7 \times 10^{-5}$ & 0.261 & 0.943 & 0.941 & 2.675 \\
        \midrule
        20  & 2 & $\sim 4.0 \times 10^{-5}$ & $\sim$0.22 & $\sim$0.95 & $\sim$0.95 & $\sim$2.70 \\
        50  & 2 & $\sim 2.5 \times 10^{-5}$ & $\sim$0.16 & $\sim$0.96 & $\sim$0.96 & $\sim$2.75 \\
        100 & 2 & $\sim 1.0 \times 10^{-6}$ & $\sim$0.12 & --- & --- & --- \\
        \bottomrule
    \end{tabular}
\end{table}

\textbf{Observation:} The validation aggregate exhibits non-monotonic behaviour within
Stage~1 (peak at epoch~4, dip at epoch~5), which is expected: stochastic gradient
descent on a randomly initialised regression head can achieve transient high validation
correlation before the head's weights overfit the training set. The best Stage~1 epoch
(epoch~4, Aggregate~$= 2.733$) is comparable to the Stage~2 early performance, indicating
that the pretrained ViT-B/32 encoder already contains substantial quality-relevant
representations without any medical domain fine-tuning.

\section{Training Loop Pseudocode}
\label{app:pseudocode}

Algorithm~\ref{alg:training} presents the pseudocode for the CT-MUSIQ two-stage
training loop.

The pseudocode for the CT-MUSIQ two-stage training loop proceeds as follows. The model
weights $\theta$ are initialised from pretrained ViT-B/32. EMA shadow weights
$\theta_{\text{EMA}}$ are set to $\theta$. An AdamW optimiser with weight decay 0.01
and a GradScaler for AMP are initialised. Then, for each epoch from 1 to 100:
(1) If epoch $\leq 5$ (Stage~1), the transformer encoder and CLS token are frozen and
learning rate is set to $3 \times 10^{-4}$. If epoch $= 6$, all parameters are
unfrozen and cosine LR scheduling is initialised.
(2) For each training batch $(\mathbf{X}, \mathbf{q})$: forward pass in fp16 autocast,
compute composite loss $\ell$, scale and backpropagate, unscale gradients, clip gradient
norm to 1.0, step optimiser, update GradScaler, and update EMA weights:
$\theta_{\text{EMA}} \leftarrow 0.999 \cdot \theta_{\text{EMA}} + 0.001 \cdot \theta$.
(3) Validation with EMA weights: temporarily substitute $\theta_{\text{EMA}}$ for $\theta$,
compute (PLCC, SROCC, KROCC, Aggregate) on the validation set, restore $\theta$.
(4) If Aggregate improves, save checkpoint and reset patience counter; otherwise increment
counter and break if counter reaches 20 (early stopping).
(5) Step cosine LR scheduler for epoch $\geq 9$.
The function returns the best saved checkpoint.

\section{Hash-Based Positional Encoding: Extended Analysis}
\label{app:hash_pos}

The hash-based positional encoding described in Section~\ref{sec:architecture} is a
key architectural contribution of CT-MUSIQ. This section provides an extended analysis
of its properties and compares it with alternative positional encoding strategies.

\subsection{Comparison with Standard Sinusoidal Encoding}

Standard sinusoidal positional encodings assign a unique, deterministic vector to each
position index $n$ using:

\begin{align}
    \text{PE}(n, 2k)   &= \sin\!\left(\frac{n}{10000^{2k/d}}\right) \\
    \text{PE}(n, 2k+1) &= \cos\!\left(\frac{n}{10000^{2k/d}}\right)
\end{align}

These encodings are well-suited for fixed-length sequences but cannot naturally encode
the two-dimensional spatial structure (row, column) or the scale identity of multi-scale
patch sequences. Naively assigning sequential position indices $0, 1, \ldots, 192$ to
the 193 tokens does not convey any information about whether a token is at scale 0 or
scale 1, or what its row/column position is within its respective grid.

\subsection{Comparison with 2D Learned Position Encodings}

Standard learned 2D positional encodings assign an independent embedding vector to each
$(row, column)$ position pair. For a fixed grid of size $G \times G$, this requires
$G^2$ embedding vectors. For our two-scale configuration:
\begin{itemize}
    \item Scale 0 ($7 \times 7$ grid): $49$ positions
    \item Scale 1 ($12 \times 12$ grid): $144$ positions
    \item Total: $49 + 144 = 193$ unique positions
\end{itemize}

However, this approach requires a separate embedding for each distinct $(scale, row,
column)$ triple, totalling $193 \times 768 = 148{,}224$ parameters. More importantly,
positions that share the same row or column index across different scales cannot share
any embedding components, preventing the model from learning that ``row~3 at scale~0''
and ``row~3 at scale~1'' occupy a similar vertical position in the image.

\subsection{Advantages of the Hash-Based Factored Encoding}

The hash-based encoding addresses both limitations through the additive factored
decomposition in \cref{eq:hash_pos}. The key advantages are:

\begin{enumerate}
    \item \textbf{Parameter efficiency}: $(2 + G_{\max} + G_{\max}) \times d
          = 28 \times 768 = 21{,}504$ parameters versus $193 \times 768 = 148{,}224$
          for the full 2D lookup table ($7\times$ reduction).

    \item \textbf{Generalisation to new scales}: If a third scale is added (e.g.,
          $512 \times 512$, $G = 16$), only $G_{\max}$ needs to be updated from 13 to 17,
          adding $(17 - 13) \times 768 \times 2 = 6{,}144$ new embedding parameters
          rather than requiring a full retraining of the positional encoding.

    \item \textbf{Row/column sharing across scales}: The row embedding at row~3 is
          shared between scale~0 and scale~1, allowing the model to learn that vertical
          position at row~3 carries similar spatial meaning regardless of scale. This
          cross-scale position sharing may improve the model's ability to correlate
          quality features across scales.

    \item \textbf{Scale embedding for scale-specific biases}: The dedicated scale
          embedding $\mathbf{E}_{\text{scale}}(s)$ allows the model to learn
          scale-specific biases (e.g., the coarse scale patches may systematically
          contribute differently to quality assessment than the fine scale patches),
          which is not possible with purely positional encodings.
\end{enumerate}

\section{Ablation Study Raw Data}
\label{app:ablation_raw}

Table~\ref{tab:ablation_raw} presents the complete validation aggregate scores for all
six ablation configurations across the 30 training epochs. Note that these are
validation aggregate values, not test set values.

\begin{table}[h]
    \centering
    \caption{Validation aggregate across ablation configurations (30 epochs).
             Values shown at selected epochs. Best configuration highlighted.}
    \label{tab:ablation_raw}
    \begin{tabular}{lcccc}
        \toprule
        \textbf{Config} & \textbf{Epoch 10} & \textbf{Epoch 20} & \textbf{Epoch 30} & \textbf{Scales} \\
        \midrule
        A1 (224, no KL)      & $\sim$2.18 & $\sim$2.23 & 2.249 & [224] \\
        A2 (224+384, no KL)  & $\sim$2.19 & $\sim$2.25 & 2.271 & [224, 384] \\
        A3 (KL=0.05)         & $\sim$2.21 & $\sim$2.26 & 2.289 & [224, 384] \\
        \textbf{A4 (KL=0.10)}& $\sim$2.23 & $\sim$2.28 & \textbf{2.305} & [224, 384] \\
        A5 (KL=0.20)         & $\sim$2.22 & $\sim$2.27 & 2.281 & [224, 384] \\
        A6 (3-scale, KL=0.10)& $\sim$2.20 & $\sim$2.26 & 2.296 & [224, 384, 512] \\
        \bottomrule
    \end{tabular}
\end{table}

\section{Model Architecture Summary Table}
\label{app:arch_summary}

Table~\ref{tab:arch_summary} provides a concise summary of the CT-MUSIQ architecture
for quick reference.

\begin{table}[h]
    \centering
    \caption{CT-MUSIQ architecture summary.}
    \label{tab:arch_summary}
    \begin{tabular}{ll}
        \toprule
        \textbf{Component} & \textbf{Specification} \\
        \midrule
        Input             & $512 \times 512$ greyscale CT slice, float32 $\in [0,1]$ \\
        Preprocessing     & Brain window (L=40, W=80 HU), channel replication $\times 3$ \\
        Scale 0           & Resize to $224 \times 224$, extract $7 \times 7 = 49$ patches \\
        Scale 1           & Resize to $384 \times 384$, extract $12 \times 12 = 144$ patches \\
        Patch size        & $32 \times 32$ pixels \\
        Total tokens      & 193 patch tokens + 1 CLS token = 194 \\
        Patch embedding   & Conv2d(3, 768, 32, 32) $\to$ Flatten $\to \mathbb{R}^{768}$ \\
        Positional encoding & $\mathbf{E}_{\text{scale}}(s) + \mathbf{E}_{\text{row}}(r) + \mathbf{E}_{\text{col}}(c)$ \\
        Transformer       & 12 layers, 8 heads, $d=768$, FFN dim$=3072$, dropout$=0.05$ \\
        Global head       & Linear(768, 384) $\to$ GELU $\to$ Dropout $\to$ Linear(384, 1) \\
        Per-scale heads   & 2 $\times$ [MeanPool $\to$ Linear(768, 384) $\to$ GELU $\to$ Linear(384, 1)] \\
        Total parameters  & $\approx 88.3$~million \\
        Output            & Scalar quality score $\hat{q} \in \mathbb{R}$ (clipped to $[0,4]$ post-hoc) \\
        \midrule
        Loss              & $\mathcal{L}_{\text{MSE}} + 0.10 \cdot \mathcal{L}_{\text{KL}} + 0.50 \cdot \mathcal{L}_{\text{rank}}$ \\
        Optimiser         & AdamW, weight decay $= 0.01$ \\
        Stage 1           & 5 epochs, LR $= 3 \times 10^{-4}$, encoder frozen \\
        Stage 2           & 95 epochs, LR: $5 \times 10^{-5} \to 10^{-6}$ (cosine) \\
        Batch size        & 10 \\
        EMA               & $\rho = 0.999$ \\
        Precision         & FP16/FP32 AMP \\
        \bottomrule
    \end{tabular}
\end{table}

\section{Dataset Statistics}
\label{app:dataset_stats}

The LDCTIQAC 2023 dataset contains 1,000 brain CT slices with radiologist Likert quality
scores. The score distribution across the full dataset has the following approximate
statistics based on the provided \texttt{train.json} labels:

\begin{itemize}
    \item \textbf{Score range}: Continuous float in $[0, 4]$
    \item \textbf{Mean score}: The majority of images receive mid-range scores (2--3),
          reflecting the clinical design of the dataset to include a range of quality levels.
    \item \textbf{Distribution}: The score distribution is approximately unimodal, with
          a peak around scores 2.5--3.0, and tails extending toward 0 (severely degraded)
          and 4 (excellent quality) that are less represented.
    \item \textbf{Training split mean}: Similar distribution to the overall dataset.
    \item \textbf{Implication for evaluation}: The use of all three correlation metrics
          (PLCC, SROCC, KROCC) rather than a single metric ensures that performance is
          evaluated across the full quality range, not just the most common quality levels.
\end{itemize}

\section{Evaluation Script Usage}
\label{app:eval_usage}

The following commands reproduce the primary evaluation results reported in this thesis:

\textbf{CT-MUSIQ evaluation (standard):}
\begin{lstlisting}[language=bash, basicstyle=\ttfamily\small, frame=single]
python evaluate.py --model ct_musiq
\end{lstlisting}

\textbf{CT-MUSIQ evaluation with TTA:}
\begin{lstlisting}[language=bash, basicstyle=\ttfamily\small, frame=single]
python evaluate.py --model ct_musiq --tta
\end{lstlisting}

\textbf{Agaldran Combo evaluation:}
\begin{lstlisting}[language=bash, basicstyle=\ttfamily\small, frame=single]
python evaluate.py --model agaldran_combo
\end{lstlisting}

\textbf{Ablation study:}
\begin{lstlisting}[language=bash, basicstyle=\ttfamily\small, frame=single]
python ablation.py
\end{lstlisting}

All results are written to CSV files in the \texttt{results/} directory and can be
compared against the published LDCTIQAC 2023 leaderboard values embedded in the
\texttt{ct\_musiq\_results.csv} and \texttt{agaldran\_combo\_results.csv} files.

\chapter{Glossary of Terms}
\label{app:glossary}

This appendix defines the key technical terms used throughout the thesis.

\begin{description}
    \item[ALARA] As Low As Reasonably Achievable. The governing principle for ionising
    radiation exposure management in clinical practice.

    \item[AMP] Automatic Mixed Precision. A training strategy that uses fp16 precision
    for forward pass activations and fp32 for gradients and parameter updates.

    \item[BRISQUE] Blind/Referenceless Image Spatial Quality Evaluator. A classical
    NR-IQA algorithm based on Natural Scene Statistics.

    \item[CLS token] Classification token. A learnable embedding vector prepended to
    the patch token sequence in Vision Transformers, used to aggregate global information
    for downstream tasks.

    \item[CT] Computed Tomography. A medical imaging modality using X-ray projections
    and computed reconstruction to produce cross-sectional anatomical images.

    \item[DLIR] Deep Learning Image Reconstruction. CT reconstruction algorithms based
    on neural networks, including CNNs and diffusion models.

    \item[EMA] Exponential Moving Average. A smoothed average of model weights used
    to reduce the stochastic variability of the final trained model.

    \item[FBP] Filtered Back Projection. The classical analytical CT reconstruction
    algorithm.

    \item[FR-IQA] Full-Reference Image Quality Assessment. Methods that evaluate image
    quality through comparison with a perfectly aligned reference image.

    \item[HU] Hounsfield Unit. The quantitative scale of CT pixel values, where water
    is 0~HU and air is approximately $-1000$~HU.

    \item[IQA] Image Quality Assessment. The task of evaluating the perceptual quality
    of an image, either through manual subjective assessment or automated algorithms.

    \item[KL divergence] Kullback--Leibler divergence. A measure of how one probability
    distribution differs from another. Used in CT-MUSIQ as a scale-consistency loss.

    \item[KROCC] Kendall Rank-Order Correlation Coefficient. A rank correlation metric
    measuring the proportion of concordant pairs among all pairs of predictions.

    \item[LDCT] Low-Dose Computed Tomography. CT acquisition with reduced radiation dose,
    achieved by lowering the X-ray tube current.

    \item[LDCTIQAC] Low-Dose Computed Tomography Perceptual Image Quality Assessment
    Grand Challenge. The MICCAI 2023 challenge that provided the benchmark dataset used
    in this thesis.

    \item[Likert scale] A rating scale used in the radiologist annotation process,
    ranging from 0 (unacceptable) to 4 (excellent) in the LDCTIQAC 2023 dataset.

    \item[MHSA] Multi-Head Self-Attention. The core computational mechanism of
    Transformer models, computing attention-weighted summaries of the input sequence.

    \item[MSE] Mean Squared Error. The primary regression loss function used in CT-MUSIQ.

    \item[MUSIQ] Multi-Scale Image Quality Transformer. The natural image IQA model
    from Ke et al.\ (ICCV 2021) that serves as the architectural foundation for CT-MUSIQ.

    \item[NPS] Noise Power Spectrum. A frequency-domain characterisation of the
    statistical properties of CT image noise.

    \item[NR-IQA] No-Reference Image Quality Assessment. Methods that evaluate image
    quality without access to any reference image.

    \item[NSS] Natural Scene Statistics. The statistical regularities exhibited by
    high-quality natural photographs, exploited by classical blind IQA algorithms.

    \item[PACS] Picture Archiving and Communication System. The hospital information
    system used to store, retrieve, and display medical images.

    \item[PLCC] Pearson Linear Correlation Coefficient. A metric measuring linear
    correlation between predicted and true quality scores.

    \item[SROCC] Spearman Rank-Order Correlation Coefficient. A metric measuring
    the agreement of rank orderings between predicted and true quality scores.

    \item[ViT] Vision Transformer. A class of image processing neural networks that
    replace convolutional layers with Transformer self-attention mechanisms.

    \item[ViT-B/32] Vision Transformer Base model with patch size 32. The pretrained
    model whose weights are used to initialise the CT-MUSIQ encoder.

    \item[VLM] Vision-Language Model. A large neural network that jointly processes
    visual and textual inputs for tasks requiring both modalities.
\end{description}
